%%%%%%%%%%%%%%%%%%%%%%%%%%%%%%%%%%%%%%%%
% RESUMO %% obrigatório
\begin{resumo}
%% neste arquivo resumo.tex o texto do resumo e as palavras-chave têm que ser em Português para os documentos escritos em Português
%% o texto do resumo e as palavras-chave têm que ser em Inglês para os documentos escritos em Inglês
%% os nomes dos comandos \begin{resumo}, \end{resumo}, \palavraschave e \palavrachave não devem ser alterados
\hypertarget{estilo:resumo}{} %% uso para este Guia


O NanoSatC-Br1 é um microssatélite brasileiro, um CubeSat executado no âmbito do Convênio do Instituto Nacional de Pesquisas Espaciais (INPE), através de sua subunidade o Centro Regional Sul de Pesquisas Espaciais (CRS) com a Universidade Federal de Santa Maria (UFSM). O OSIRIS-REx é uma missão de ciência planetária, que consiste em estudar e coletar amostras do asteroide 101955 Bennu, um asteroide carbonáceo, trazendo a amostra até a Terra em 2023. O objetivo deste trabalho é mostrar os aspectos principais de cada missão, como objetivo, descrição, funcionamento, e equipamentos científicos.




\palavraschave{
	\palavrachave{Missões, Espaço}
}


\end{resumo}